\documentclass{article}
\usepackage{amsmath}
\usepackage{amsthm}

\newtheorem{theorem}{Problem}

\begin{document}

\begin{theorem}
How many positive two-digit integers are factors of both 100 and 150?
\end{theorem}

\begin{proof}
Step 0: To find the factors of 100 and 150, I can use prime factorization and list all the possible combinations of the prime factors.
Step 1: The prime factorization of 100 is 2^2 times 5^2, so the factors of 100 are 1, 2, 4, 5, 10, 20, 25, 50, and 100.
Step 2: The prime factorization of 150 is 2 times 3 times 5^2, so the factors of 150 are 1, 2, 3, 5, 6, 10, 15, 25, 30, 50, 75, and 150.
Step 3: To find the common factors of both 100 and 150, I can look for the numbers that appear in both lists.
Step 4: These are 1, 2, 5, 10, and 50.
Step 5: However, the question asks for the positive two-digit integers, so I can eliminate 1, 2, and 5 from the list of common factors.
Step 6: Therefore, the only positive two-digit integers that are factors of both 100 and 150 are 10 and 50, and there are two of them.
Step 7: # Answer\n2
\end{proof}

\end{document}
